% book/chapters/ch05-skills.tex
\chapter{Skills}
\label{ch:skills}

\index{skills}

A skill is a reusable instruction set stored in a \texttt{SKILL.md} file.
Where a subagent defines \emph{who} does a task (specialised model, restricted
tools), a skill defines \emph{how} a task is done — the checklist, workflow,
or procedure Claude should follow.

Claude Code skills follow the \textbf{Agent Skills open standard}
(\url{https://agentskills.io}), which means skill files created here can work
across multiple AI tools that implement the same standard.

\section{Skills vs.\ CLAUDE.md vs.\ Subagents}
\label{sec:skill-comparison}

\begin{table}[ht]
\centering
\caption{Choosing between CLAUDE.md, Skills, and Subagents}
\label{tab:skill-vs}
\begin{tabularx}{\textwidth}{lXX}
\toprule
\textbf{Mechanism} & \textbf{Best for} & \textbf{Loads} \\
\midrule
\texttt{CLAUDE.md}    & Persistent facts and conventions every session needs
                      & Every session, always \\
Skill                 & Procedures and workflows triggered only sometimes
                      & On demand (no cost until invoked) \\
Subagent              & Isolated task execution with its own context window
                      & On demand \\
\bottomrule
\end{tabularx}
\end{table}

The key advantage of skills over \texttt{CLAUDE.md} content: skill bodies load
\emph{only when invoked}, so long reference material costs almost nothing until
needed.

\section{Skill Storage and Scopes}
\label{sec:skill-scopes}

\index{skills!scopes}

\begin{table}[ht]
\centering
\caption{Skill storage locations}
\label{tab:skill-scopes}
\begin{tabularx}{\textwidth}{llX}
\toprule
\textbf{Scope} & \textbf{Path} & \textbf{Applies to} \\
\midrule
Enterprise & Managed settings & All users in the organisation \\
Personal   & \texttt{\textasciitilde/.claude/skills/<name>/SKILL.md}
           & All your projects \\
Project    & \texttt{.claude/skills/<name>/SKILL.md}
           & This project only \\
Plugin     & \texttt{<plugin>/skills/<name>/SKILL.md}
           & Where plugin is enabled (namespaced) \\
\bottomrule
\end{tabularx}
\end{table}

When skill names conflict, enterprise overrides personal, which overrides
project. Plugin skills are namespaced (\texttt{plugin:skill}) and never
conflict.

\begin{notebox}
Files in \texttt{.claude/commands/} still work and create the same
\texttt{/command} shortcuts. Skills supersede commands when both share a name,
and skills support additional features (supporting files, invocation control,
subagent execution).
\end{notebox}

\section{Creating a Skill}
\label{sec:skill-create}

\index{skills!creating}
\index{SKILL.md}

Each skill is a directory containing a \texttt{SKILL.md} file plus optional
supporting files:

\begin{lstlisting}[language=bash,caption={Skill directory structure}]
.claude/skills/fix-issue/
  SKILL.md          # main instructions (required)
  examples/
    sample-output.md
  scripts/
    run-tests.sh
\end{lstlisting}

\subsection{SKILL.md Frontmatter}

\begin{lstlisting}[caption={Example \texttt{SKILL.md} with frontmatter}]
---
name: fix-issue
description: Fixes a GitHub issue end-to-end. Claude invokes automatically
             when asked to fix or resolve a GitHub issue.
model: sonnet
disable-model-invocation: false
---

Analyse and fix the GitHub issue: $ARGUMENTS

1. Run `gh issue view $ARGUMENTS` to read the issue
2. Search the codebase for relevant files
3. Implement the fix
4. Write a failing test that reproduces the issue, then verify it passes
5. Ensure linting and type checks pass
6. Create a descriptive commit and open a PR with `gh pr create`
\end{lstlisting}

\begin{table}[ht]
\centering
\caption{SKILL.md frontmatter fields}
\label{tab:skill-frontmatter}
\begin{tabularx}{\textwidth}{llX}
\toprule
\textbf{Field} & \textbf{Required} & \textbf{Description} \\
\midrule
\texttt{name}                     & Yes & Identifier; becomes the \texttt{/name} slash command \\
\texttt{description}              & Yes & When Claude auto-invokes this skill \\
\texttt{model}                    & No  & \texttt{haiku}, \texttt{sonnet}, or \texttt{opus} \\
\texttt{disable-model-invocation} & No  & \texttt{true} requires explicit slash-command invocation \\
\bottomrule
\end{tabularx}
\end{table}

\subsection{The \$ARGUMENTS Variable}

\texttt{\$ARGUMENTS} captures everything typed after the skill name when invoked
with a slash command:

\begin{lstlisting}[caption={Invoking a skill with arguments}]
/fix-issue 1234
# $ARGUMENTS = "1234"

/code-review --deep src/auth.ts
# $ARGUMENTS = "--deep src/auth.ts"
\end{lstlisting}

\section{Invoking Skills}
\label{sec:skill-invoke}

\index{skills!invocation}

Skills can be invoked in three ways:

\begin{enumerate}
  \item \textbf{Slash command}: type \texttt{/skill-name} (with optional
        arguments) in any session.
  \item \textbf{Natural language}: Claude auto-invokes the skill when the
        request matches the \texttt{description} field, unless
        \texttt{disable-model-invocation: true} is set.
  \item \textbf{Programmatic}: \texttt{claude -p "/fix-issue 1234"} in
        non-interactive mode.
\end{enumerate}

\section{Running Skills in a Subagent}
\label{sec:skill-subagent}

\index{skills!subagent execution}

A skill can be executed inside a subagent to isolate its context from the main
conversation. Specify this in the frontmatter:

\begin{lstlisting}[caption={Running a skill in a subagent (frontmatter option)}]
---
name: heavy-analysis
description: Deep analysis that reads many files
run-in-subagent: true
model: haiku
---
\end{lstlisting}

\section{Installing Marketplace Skills}
\label{sec:skill-marketplace}

\index{skills!marketplace}

The superpowers collection (\url{https://github.com/obra/superpowers}) provides
production-ready skills for common workflows:

\begin{lstlisting}[language=bash,caption={Installing the official superpowers plugin}]
/plugin install superpowers@claude-plugins-official
/reload-plugins
\end{lstlisting}

After installation, skills such as \texttt{/code-review}, \texttt{/test-driven-development},
and \texttt{/writing-plans} are available immediately.

%% ---------------------------------------------------------------
\begin{summarybox}
\begin{itemize}
  \item A skill is a \texttt{SKILL.md} file that defines \emph{how} a task
        is done; it loads only when invoked, unlike \texttt{CLAUDE.md} which
        loads every session.
  \item Skills follow the Agent Skills open standard (\url{https://agentskills.io})
        and can work across multiple AI tools.
  \item Skills are stored in \texttt{.claude/skills/<name>/SKILL.md} (project)
        or \texttt{\textasciitilde/.claude/skills/<name>/SKILL.md} (personal).
  \item The \texttt{name} field becomes a \texttt{/name} slash command;
        the \texttt{description} field controls auto-invocation; \texttt{model}
        routes to a specific Claude model; \texttt{\$ARGUMENTS} captures
        command-line input.
  \item Custom commands (\texttt{.claude/commands/}) are still supported and
        behave identically to skills.
\end{itemize}
\end{summarybox}

\printbibliography[heading=subbibliography,title={References}]
